% US-082 DE4 (Lower Saxony) transfer experiment --- the positive counterpart of the
% multi-region study. \input-able section: no \documentclass / \begin{document}.
% Prose in English. Every figure is a real EPIC 12 artefact; every number is anchored
% with a % src: comment to its source (report.json / engram-confirmed run).
% This section is \input by main.tex right after experiments_multiregion.tex.

\subsection{Una región donde la transferencia funciona: Baja Sajonia (DE4)}
\label{sec:de4}

El estudio multirregión anterior es, por diseño, un catálogo de
\emph{negativos honestos}: el zero-shot falla en las brechas
franco-ibérica y tropical, y agrupar regiones amplía la taxonomía sin
rescatar una sola clase fina por encima de la calidad de producción.
Surge una pregunta natural: \emph{¿existe una región objetivo donde el
campeón entrenado en Francia transfiera bien, y qué la hace diferente?}
Esta subsección responde que sí, y aísla las tres palancas que lo
deciden. El objetivo es Baja Sajonia, Alemania (región NUTS DE4),
elegida porque su cobertura en EuroCrops~v2 es \emph{completa} (no
parcial como en el piloto de la Toscana italiana) y su agricultura está
dominada por cereales de invierno grandes y fenológicamente separables.

\paragraph{Cómo se alcanzó el techo de Toscana: tres estrategias de transfer.}
Antes de DE4, el mismo backbone entrenado en Francia (TSViT-pheno) se llevó al
dataset italiano de Toscana ($39$ clases nativas, test fold-3) bajo tres
estrategias de transfer, para separar qué aporta realmente la transferencia.
\emph{(i)} Entrenar desde cero solo con datos italianos alcanza F1-macro~$0.095$.
\emph{(ii)} El transfer puro --- reutilizar el codificador temporal de Francia y
aprender un clasificador nuevo para los cultivos italianos --- lo eleva a $0.120$.
\emph{(iii)} Conservar las clases francesas (arrancar la cabeza desde los cultivos
franceses) o replicar el modo de entrenamiento completo de PASTIS llega a
$0.130$--$0.132$. El codificador temporal es, por tanto, el componente transferido
determinante; conservar las clases suma poco en promedio ($+0.012$) y \emph{no} es
universal --- ayuda a cultivos parecidos entre Francia e Italia (trigo duro $+0.08$,
cebada $+0.07$, alfalfa $+0.05$) y perjudica a los de fenología mediterránea distinta
(avena $-0.08$, girasol $-0.06$). Este techo de Toscana de $\approx0.13$ en el mejor
caso es la línea base que las palancas de DE4 a continuación elevan a $0.266$.
% src: reports/us079_figs/ablation_compare.json (A_warmstart 0.132, B_nowarmstart 0.120,
% original 0.130; from-scratch 0.095; deltas warm-start por clase)

\paragraph{Construir un dataset homólogo a PASTIS para DE4.}
Materializamos un dataset denso homólogo a PASTIS para DE4 2023 desde
cero, con \textbf{100\% de datos reales y cero marcadores de posición}:
series temporales Sentinel-2 L2A por parche
(\texttt{DATA\_S2/S2\_*.npy}, 10 bandas, $128\times128$) obtenidas del
Copernicus Data Space Ecosystem (CDSE) Process API, máscaras densas de
cultivo (\texttt{ANNOTATIONS/TARGET\_*.npy}) rasterizadas a partir de los
polígonos oficiales de EuroCrops~v2 (JRC, taxonomía HCAT4), y el
embedding anual de AlphaEarth de 64 dimensiones
(\texttt{SATELLITE\_EMBEDDING/V1/ANNUAL} v1.1, año 2023) obtenido vía
Google Earth Engine para cada parcela etiquetada. El espacio de
etiquetas es el conjunto nativo de cultivos HCAT alemán (37 clases),
construido por el mismo generador region-generalizado usado para Italia
(\texttt{--region-prefix de4}); cada código HCAT se verifica contra el
crosswalk oficial del JRC, ninguno se inventa.

\paragraph{Tres palancas, medidas de forma independiente.}
Partiendo del techo del piloto de la Toscana (Voting-3 F1
macro~$\approx0.12$), identificamos y medimos tres palancas
independientes, cada una sobre datos out-of-fold reales
(Tabla~\ref{tab:de4_levers}):
\begin{enumerate}[leftmargin=1.4em,itemsep=2pt]
  \item \textbf{Año de adquisición 2023 vs.\ 2018}: en la Toscana, el
    miembro XGBoost-sobre-AlphaEarth sube de F1~$0.156$ a $0.225$
    ($+44\%$) puramente al mover el embedding anual a un año más reciente
    y mejor cubierto. % src: us-082-REPORTE-RESULTADOS.md Table 2
  \item \textbf{Ventana temporal de 14 meses vs.\ 8}: PASTIS-R abarca 13
    meses (sep~2018--oct~2019, incluyendo el invierno), mientras que el
    piloto italiano cubría solo mar--oct (8 meses), privando de señal a
    la fenología de los cereales de invierno. Reextraer con una ventana
    \texttt{2022-11-01..2023-12-31} (y un \texttt{max\_cloud}~$=50$
    relajado) eleva el número de fechas utilizables de 24 a 56 en la
    Toscana y a $\sim$41 en DE4 --- cerca de las 43 fechas de PASTIS.
    % src: engram HALLAZGO VENTANA TEMPORAL US-082
  \item \textbf{Región DE4 vs.\ Toscana}: la palanca dominante. Con
    parcelas grandes y limpias (mediana de 10\,ha vs.\ 0.4\,ha) y
    cobertura de cultivos completa, el F1 macro de Voting-3 salta de
    $0.119$ a $0.266$ ($2.2\times$), y los cereales de invierno --- trigo,
    maíz, centeno, colza --- se recuperan de F1~$0.05$--$0.20$ a
    $0.57$--$0.70$. % src: report.json de4_2023
\end{enumerate}

\begin{table}[t]
  \centering
  \caption{Las tres palancas independientes que elevan la transferencia
  desde el techo de la Toscana hasta DE4, cada una medida sobre datos
  out-of-fold reales. El cambio de región (Toscana$\rightarrow$DE4) es el
  factor dominante: los campos grandes y separables de cereal de invierno
  cierran la brecha que las parcelas mediterráneas pequeñas y fragmentadas
  no podían. Miembro XGBoost-sobre-AlphaEarth salvo que se indique.}
  \label{tab:de4_levers}
  \begin{tabular}{lcc}
    \toprule
    Configuración & XGB F1 macro & Exactitud \\
    \midrule
    Toscana 2018 (ventana de 8 meses, 24 fechas)   & 0.156 & 0.37 \\
    Toscana 2023 (ventana de 14 meses, 56 fechas)  & 0.225 & 0.45 \\
    DE4 2023 (ventana de 14 meses, $\sim$41 fechas) & \textbf{0.267} & \textbf{0.633} \\
    \bottomrule
  \end{tabular}
\end{table}

\paragraph{El campeón Voting-3 sobre DE4.}
Reejecutamos la receta completa del campeón sobre DE4 --- TSViT-pheno
(capacidad Full-M, inicializado en caliente desde el campeón de PASTIS)
más U-TAE más XGBoost-sobre-AlphaEarth, agregados a parcelas y
combinados por el ensamble de votación de tres miembros
(\texttt{ItaliaParcelVotingEnsemble}). El mejor miembro denso,
TSViT-pheno-fullm, alcanza F1~$0.333$ a nivel fino y exactitud~$0.682$ en
el fold de prueba retenido, y el Voting-3 obtiene un F1
macro~$0.266$ % src: report.json de4_2023
--- una mejora de $2.2\times$ sobre el Voting-3 de la Toscana ($0.119$).
Crucialmente, la anatomía por clase confirma la tesis de cardinalidad de
la Sección~\ref{sec:multiregion:perclass} desde el lado \emph{positivo}:
donde DE4 tiene campos grandes y separables, los cereales previamente
colapsados se rescatan (maíz~$0.703$, colza~de~invierno~$0.689$,
trigo~de~invierno~$0.674$, centeno~de~invierno~$0.573$), arrojando
8~clases con F1~$\ge0.60$ y 12 con $\ge0.40$, frente a una única clase
por encima de $0.60$ en la Toscana. % src: report.json de4_2023 per-class

\paragraph{Tres convenciones de espacio de etiquetas.}
Para una lectura honesta evaluamos el mismo ensamble entrenado bajo tres
convenciones de espacio de etiquetas, decididas a~priori: \emph{(A)
conservando las 37 clases HCAT nativas} (fino, F1 macro~$0.266$);
\emph{(B) sin conservar}, proyectando la predicción sobre el espacio
grueso al estilo PASTIS mediante crosswalk (grueso, F1 macro~$0.287$); y
\emph{(C) el procedimiento completo de extremo a extremo} (la propia
votación entrenada, $0.266$). % src: report.json de4_2023 fine/coarse
La vista gruesa (B) es la más favorable, como cabía esperar cuando se
fusionan cereales hermanos casi inseparables, pero reportamos las tres
para que la cifra no se seleccione tendenciosamente por taxonomía.

\paragraph{Una rama fenológica medida para DE4.}
Un punto metodológico sutil pero importante: la rama semántica de
TSViT-pheno alinea cada píxel contra un prototipo fenológico por clase
--- una descripción de curva NDVI media generada por Gemini y codificada
por un sentence transformer (siguiendo a Wen et al.,
2025~\cite{wen2025phenology}). Los prototipos son específicos de la
región: un calendario alemán de cereal de invierno, no el mediterráneo.
Generamos los prototipos de DE4 a partir de las curvas NDVI reales por
clase de los parches de DE4, con la descripción anclada en el calendario
de Baja Sajonia 2023; las curvas son reales y solo el calificador
textual es regional, nunca fabricado. Es la misma disciplina aplicada en
todo el trabajo: señales reales, anclaje regional, sin datos sintéticos.

\paragraph{Salvedades honestas.}
El Voting-3 de DE4 se mide sobre un único fold de prueba retenido; la
cifra de titular inatacable requeriría un protocolo out-of-fold de 5
folds ($5\times$ de cómputo), aunque el patrón por clase es decisivo y no
cambiaría el veredicto. La cobertura espacial de la muestra de DE4 es
parcial --- las parcelas alemanas son grandes y dispersas, de modo que
los parches etiquetados cubren una fracción del universo regional ---
pero la pérdida por clase es uniforme entre clases (sin sesgo de
selección), de modo que la muestra es representativa. DE4 ya gana de
forma decisiva con esta cobertura; ampliarla solo mejora la estimación.
Este es el broche positivo del estudio multirregión: la misma receta que
falla en zero-shot a través de brechas distantes \emph{tiene éxito}
cuando la región objetivo aporta campos grandes y fenológicamente
separables y un registro temporal completo, reciente e inclusivo del
invierno.
